% ===== 原文 pp.147–159（PDF p.162–174）+ p.160 页首定理收尾，tr-03：第 4 章 §4.3 后段与 §4.4 前段 =====
% 块首不写 §4.3 标题（属 c04a，位于印刷 p.143）；§4.4 标题按原文位置写在下方。

% [原文 p.147]
\begin{definition}
$X_0$ 中的\emph{环面横截前对数曲线}是指一个稳定映射 $f:C\to X_0$，其中 $C$ 是曲线，使得对每个顶点 $v\in\cP$，$f^{-1}(D_v)\to D_v$ 是环面横截曲线，并且 $f$ 满足命题 4.9 的结论 (1) 与 (2)。
\end{definition}

由于我们只用了命题 4.9 的 (1) 与 (2) 来构造与环面横截对数曲线相伴的热带曲线，事实上我们可以对任何环面横截前对数曲线相伴一条热带曲线。

现在回到本节开头的情形：给定一条标记热带曲线
\[
h:(\Gamma,x_1,\dots,x_s)\to M_\R
\]
满足 $h(x_i)=P_i$，并与之相伴地给出带有额外顶点的曲线 $h:\tilde\Gamma\to M_\R$。我们可以对 $h:\Gamma\to M_\R$ 相伴另一条热带曲线 $h:\hat\Gamma\to M_\R$，它由从 $\Gamma$ 中删去所有标记边、并照常删去由此产生的二价顶点而得到。于是图 $\hat\Gamma$ 有可区分的边 $E_1,\dots,E_s$，使得 $E_{x_i}$ 在 $\Gamma$ 中的端点现在落在 $E_i$ 的内部。先验地注意，诸 $E_i$ 不必互不相同（事实上它们会互不相同，但我们用不到这一点）。若 $\hat\Gamma$ 的一条边 $E$ 对某个 $i$ 等于 $E_i$，我们就说 $E$ 是 $\hat\Gamma$ 的一条\emph{标记边}。

\begin{proposition}
对 $\hat\Gamma$ 的每条边 $E$ 选取一个定向。这给出 $E$ 的端点的一种选取，我们记之为 $\partial^-E$ 与 $\partial^+E$（若 $E$ 非紧，则只记 $\partial^-E$）。做完这一选取之后，令
\[
u_{(\partial^-E,E)}\in M
\]
表示 $h(E)$ 的一个本原切向量，它从 $h(\partial^-E)$ 指向 $h(\partial^+E)$（或指向 $h(E)$ 的无界方向）。则映射
\begin{equation}
\begin{aligned}
\Phi&:\operatorname{Map}(\hat\Gamma^{[0]},M)\to\Big(\prod_{E\in\hat\Gamma^{[1]}\setminus\hat\Gamma^{[1]}_\infty}M/\Z u_{(\partial^-E,E)}\Big)\times\Big(\prod_{i=1}^sM/\Z u_{(\partial^-E_i,E_i)}\Big)\\
H&\mapsto\Big((H(\partial^+E)-H(\partial^-E))_E,(H(\partial^-E_i))_i\Big).
\end{aligned}
\tag{4.6}
\end{equation}
是有限指数的格包含。称这个指数为 $\mathfrak{D}$。
\end{proposition}

于是 $\mathfrak{D}$ 就是形如 $f:(C,x_1,\dots,x_s)\to X_0$、满足 $f(x_i)=q_i$ 且相伴热带曲线为 $h$ 的标记环面横截前对数曲线在同构意义下的个数。

% [Tutor Note: p.147 的 $\mathfrak{D}$ 是映射 $\Phi$ 的像在格中的指数。命题 4.22 末句把 $\mathfrak{D}$ 解释为前对数曲线的个数，核心是：$\Phi\otimes\R$ 的核对应热带曲线的形变自由度（假设为零），而 $\Phi$ 的余核阶（等于像的指数）对应粘合的自由度。]

\begin{proof}[证明]
把 $\Phi$ 与 $\R$ 作张量积，我们便得到保有点 $P_i$ 的关联条件的、热带曲线 $h:\hat\Gamma\to M_\R$ 的形变的一种描述。确实，给定 $H\in\operatorname{Map}(\hat\Gamma^{[0]},M_\R)$，我们可以通过把顶点 $V$ 送到 $h(V)+H(V)$ 来修改 $h$。此时紧边 $E$ 被映到连接 $h(\partial^-E)+H(\partial^-E)$ 与 $h(\partial^+E)+H(\partial^+E)$ 的边。非紧且非收缩的边 $E$ 被映到射线 $\R_{\ge0}u_{(\partial^-E,E)}+h(\partial^-E)+H(\partial^-E)$。一般地，这样的形变并非热带曲线；而若我们想把 $h$ 形变为仍为热带曲线的曲线，不改变组合类型，并保证 $E_i$ 的像仍包含 $P_i$，则需要两个条件成立：

\begin{enumerate}[label=(\arabic*)]
\item 对每条紧边 $E$，$(h(\partial^+E)+H(\partial^+E))-(h(\partial^-E)+H(\partial^-E))$ 平行于 $u_{(\partial^-E,E)}$；等价地，
\[
H(\partial^+E)-H(\partial^-E)=0\quad\text{于 }M/\Z u_{(\partial^-E,E)}.
\]
\end{enumerate}

% [原文 p.148]
\begin{enumerate}[label=(\arabic*),start=2]
\item 对每个 $i$，$H(\partial^-E_i)=0$ 于 $M/\Z u_{(\partial^-E_i,E_i)}$。这一点连同 $E=E_i$ 时的 (1)，保证了 $E_i$ 的像所张成的仿射直线仍经过 $P_i$。
\end{enumerate}

更简洁地说，$H$ 给出 $h$ 作为热带曲线的一个形变，仅当 $H$ 属于 $\Phi\otimes\R$ 的核。若 $H$ 离原点足够近，则它会给出 $h$ 的一个形变。由于按假设 $h$ 是刚性的，$\Phi$ 的核事实上是零。于是维数计数表明，$\Phi$ 的像是有限指数的一个子格，记为 $\mathfrak{D}$。

现在对每个 $V\in\tilde\Gamma^{[0]}$，记 $\mathbf{u}(V),\mathbf{w}(V)$ 为与 $V$ 相邻的未标记边的本原切向量（从 $h(V)$ 指向外）及其权所构成的数据。

我们所要构造的曲线 $C$ 必然是由环面曲面 $D_{h(V)}$ 中的直线粘合而成，其中 $V$ 跑遍 $\tilde\Gamma^{[0]}$。对给定的 $V$，$D_{h(V)}$ 中类型为 $(\mathbf{u}(v),\mathbf{w}(v))$ 的直线之集是 $\cL_{(\mathbf{u}(V),\mathbf{w}(V))}$。现在，对无相邻标记边的三价顶点 $V\in\tilde\Gamma^{[0]}$，由推论 4.20，$\cL_{(\mathbf{u}(V),\mathbf{w}(V))}$ 是 $\mathbb{G}(M)$ 上的 torsor。对二价的、或有相邻标记边的三价顶点 $V\in\tilde\Gamma^{[0]}$，设其相邻未标记边为 $E^\pm(V)$，则由推论 4.18，$\cL_{(\mathbf{u}(V),\mathbf{w}(V))}$ 是 $\mathbb{G}(M/u_{(V,E^\pm(V))})$ 上的 torsor。注意无相邻标记边的三价顶点恰是 $\hat\Gamma$ 的顶点，而其余顶点正是 $\tilde\Gamma$ 的那些不是 $\hat\Gamma$ 顶点的顶点。于是
\[
\prod_{V\in\tilde\Gamma^{[0]}}\cL_{(\mathbf{u}(V),\mathbf{w}(V))}
\]
是
\[
\mathbb{G}(\operatorname{Map}(\hat\Gamma^{[0]},M))\times\prod_{V\in\tilde\Gamma^{[0]}\setminus\hat\Gamma^{[0]}}\mathbb{G}(M/u_{(V,E^-(V))})
\]
上的 torsor。

为了得到一条环面横截前对数曲线，我们需要把各种直线与标记点匹配起来。我们通过定义映射
\begin{equation}
\Phi':\prod_{V\in\tilde\Gamma^{[0]}}\cL_{(\mathbf{u}(V),\mathbf{w}(V))}\to\prod_{E\in\tilde\Gamma^{[1]}\setminus\tilde\Gamma^{[1]}_\infty}\mathbb{G}(M/u_{(\partial^-E,E)})\times\prod_{i=1}^s\mathbb{G}(M/\Z u_{(\partial^-E_i,E_i)})
\tag{4.7}
\end{equation}
来做这件事，具体如下。

% [Tutor Note（原文笔误，照译未改）: 式 (4.7) 中第一个因子写作 $\mathbb{G}(M/u_{(\partial^-E,E)})$，少了一个 $\Z$；同页下文说该轨道同构于 $\mathbb{G}(M/\Z u_{(\partial^-E,E)})$，而式 (4.6)、(4.8) 均带 $\Z$。译文按原文作 $M/u$。]

(1) 给定该映射定义域中的一个元素 $(\varphi_V)_{V\in\tilde\Gamma^{[0]}}$，则它在 $\Phi'$ 下的像中对应于 $E\in\tilde\Gamma^{[1]}\setminus\tilde\Gamma^{[1]}_\infty$ 的分量定义如下。令 $V^\pm=\partial^\pm E$，$v^\pm=h(V^\pm)$，并令 $\omega$ 为 $\cP$ 中连接 $v^+$ 与 $v^-$ 的边，于是 $D_{v^+}\cap D_{v^-}=D_\omega$。我们有 $\varphi_{V^\pm}:C_\pm\cong\PP^1\to D_{v^\pm}$，设 $p^\pm=\varphi_{v^\pm}(C_\pm)\cap D_\omega$。由横截性，两点 $p^\pm$ 都落在 $D_\omega$ 的大环面轨道中，而该轨道同构于 $\mathbb{G}(M/\Z u_{(\partial^-E,E)})$。我们便取 $(\varphi_V)_V$ 在 $\Phi'$ 下的像中对应于 $E$ 的分量为比值 $p^+/p^-$。该比值为 $1$ 恰意味着 $p^-=p^+$。

(2) 给定一个元素 $(\varphi_V)_{V\in\tilde\Gamma^{[0]}}$，则 $\Phi'$ 的像中对应于 $E_i$ 的分量定义如下。令 $V_i\in\tilde\Gamma^{[0]}$ 为与 $E_{x_i}$ 相邻的顶点，于是 $h(V_i)=P_i$。由于 $P_i$ 落在 $h(E_i)$ 的内部，投影 $M\to M/\Z u_{(\partial^-E_i,E_i)}$ 定义了环面簇之间的映射 $g_i:D_{P_i}\to\PP^1\supseteq\mathbb{G}(M/\Z u_{(\partial^-E_i,E_i)})$。由引理 4.17，$\varphi_{V_i}$ 的像是纤维 $g_i^{-1}(r_i)$。我们便取 $(\varphi_V)_V$ 在 $\Phi'$ 下的像中对应于 $E_i$ 的分量为 $(g_i(q_i))/r_i$。这是有意义的。确实，由命题 4.13，$q_i\in D_{P_i}\setminus(\partial D_{P_i})$，故 $g_i(q_i)\in\PP^1\setminus\partial\PP^1$。注意该比值为 $1$ 恰当 $q_i$ 落在 $\varphi_{V_i}$ 的像中。

% [原文 p.149]
我们可以这样使用映射 $\Phi'$。给定一条环面横截前对数曲线 $\varphi:(C,x_1,\dots,x_s)\to X$，满足 $\varphi(x_i)=q_i$ 且相伴热带曲线为 $h$，则 $\tilde\Gamma^{[0]}$ 的元素与 $C$ 的不可约分量的等价类一一对应。对 $V\in\tilde\Gamma^{[0]}$，若 $C_V$ 表示对应于 $V$ 的等价类中各不可约分量的并，则由于 $h$ 是单的，$\varphi_V:=\varphi|_{C_V}:C_V\to D_{h(V)}$ 是一条直线。确实，由于 $h$ 是单的，$C_V$ 上至多有 $3$ 个点映入 $\operatorname{Sing}(X_0)$，因为按构造这样的点与 $\Gamma$ 中与 $V$ 相邻的有界边一一对应。此外仍由单性，这些点各自映到 $\operatorname{Sing}(X_0)$ 的不同不可约分量，因为 $\Gamma$ 中与 $V$ 相邻的两条边不可能映到 $M_\R$ 中的同一线段。由于 $C_V$ 的任一非收缩分量的像与 $\operatorname{Sing}(X_0)$ 至少交于两点，$C_V$ 至多含一个非收缩分量。又因 $C$ 是有理曲线且 $\varphi$ 是稳定的，事实上 $C_V$ 不能含任何收缩分量。因此 $C_V$ 不可约，且 $\varphi_V$ 是一条直线。我们得到
\[
(\varphi_V)_V\in\prod_{V\in\tilde\Gamma^{[0]}}\cL_{(\mathbf{u}(V),\mathbf{w}(V))}.
\]
此外，该元素在 $\Phi'$ 下的像为 $(1,\dots,1)$。确实，上面类型 (1) 的分量为 $1$ 恰是说直线 $\varphi_{\partial^-E}$ 与 $\varphi_{\partial^+E}$ 经过 $D_\omega$ 中的同一点；而上面类型 (2) 的分量为 $1$ 是说直线 $\varphi_{V_i}$ 经过点 $q_i$。

反之，给定 $(\varphi_V)_V\in\prod_{V\in\tilde\Gamma^{[0]}}\cL_{(\mathbf{u}(V),\mathbf{w}(V))}$ 使它在 $\Phi'$ 下的像为 $(1,\dots,1)$，显然我们可以倒转该过程，把诸 $C_V$ 粘合起来，得到一条环面横截前对数曲线 $\varphi:C\to X_0$。我们还需要标记这条曲线：对 $1\le i\le s$，选取一点 $x_i\in C_{V_i}$，使它在 $\varphi_{V_i}$ 下映到 $q_i$。这有 $w(E_i)$ 种选择；但 $C_{V_i}$ 有保持 $\varphi_{V_i}^{-1}(\partial D_{P_i})$ 的两个点不动的自同构，它把 $x_i$ 的任一种选取送到另一种选取。于是我们断言：在同构意义下，经过 $q_1,\dots,q_s$ 且相伴热带曲线为 $h$ 的环面横截前对数映射 $f:(C,x_1,\dots,x_s)\to X_0$ 之集，恰是 $(1,\dots,1)$ 在 $\Phi'$ 下的原像。我们希望计算该集合的基数，即 $\Phi'$ 的次数。

现在 $\Phi'$ 是代数环面上的 torsor 之间的映射，其对应的特征格映射为 $\Phi''$，后者可如下描述。对每个顶点 $V\in\tilde\Gamma^{[0]}\setminus\hat\Gamma^{[0]}$，令 $E(V)$ 记包含 $V$ 的那条唯一的 $\hat\Gamma$ 的边。对 $\tilde\Gamma$ 的一条边 $E$，记 $\hat E$ 为包含 $E$ 的 $\hat\Gamma$ 的边。于是我们可以写
\begin{equation}
\begin{aligned}
\Phi'':{}&\prod_{V\in\hat\Gamma^{[0]}}M\times\prod_{V\in\tilde\Gamma^{[0]}\setminus\hat\Gamma^{[0]}}M/\Z u_{(\partial^-E(V),E(V))}\\
&\to\prod_{E\in\tilde\Gamma^{[1]}\setminus\tilde\Gamma^{[1]}_\infty}M/\Z u_{(\partial^-\hat E,\hat E)}\times\prod_{i=1}^sM/\Z u_{(\partial^-E_i,E_i)}.
\end{aligned}
\tag{4.8}
\end{equation}
给定 $(m_V)_{V\in\tilde\Gamma^{[0]}}$，其中若 $V\in\hat\Gamma^{[0]}$ 则 $m_V\in M$，而若 $V\in\tilde\Gamma^{[0]}\setminus\hat\Gamma^{[0]}$ 则 $m_V\in M/\Z u_{(\partial^-E(V),E(V))}$，则 $(m_V)$ 在 $\Phi''$ 下的像有如下分量。对应于边 $E\in\tilde\Gamma^{[1]}\setminus\tilde\Gamma^{[1]}_\infty$ 的分量为 $m_{\partial^+E}-m_{\partial^-E}$。对应于 $E_i$ 的分量为 $m_{V_i}$。

% [Tutor Note: 本块中 $\hat\Gamma$ 与 $\tilde\Gamma$ 是两种不同的图：$\tilde\Gamma$ 由 $\Gamma$ 在每条标记边内部添加一个额外顶点得到，$\hat\Gamma$ 则由删去全部标记边得到。式 (4.6) 的乘积跑遍 $\hat\Gamma$ 的边，式 (4.7) 的乘积跑遍 $\tilde\Gamma$ 的边，式 (4.8) 的源按 $\hat\Gamma^{[0]}$ 与 $\tilde\Gamma^{[0]}\setminus\hat\Gamma^{[0]}$ 拆分；核对时勿混。]

我们可以把 $\Phi''$ 与 $\Phi$ 联系起来，如下。令 $S\subseteq\tilde\Gamma^{[1]}\setminus\tilde\Gamma^{[1]}_\infty$ 为使 $\hat E$ 紧且 $\partial^+\hat E$ 是 $E$ 的一个顶点的那些边 $E$ 之集；显然 $S$ 与 $\hat\Gamma^{[1]}\setminus\hat\Gamma^{[1]}_\infty$ 之间有双射。

% [原文 p.150]
于是我们拆分
\[
\prod_{E\in\tilde\Gamma^{[1]}\setminus\tilde\Gamma^{[1]}_\infty}M/\Z u_{(\partial^-\hat E,\hat E)}=\prod_{E\in S}M/\Z u_{(\partial^-\hat E,\hat E)}\times\prod_{E\in\tilde\Gamma^{[1]}\setminus(\tilde\Gamma^{[1]}_\infty\cup S)}M/\Z u_{(\partial^-\hat E,\hat E)}.
\]
现在 $\Phi''$ 在
\[
\prod_{V\in\tilde\Gamma^{[0]}\setminus\hat\Gamma^{[0]}}M/\Z u_{(\partial^-E(V),E(V))}
\]
与
\[
\prod_{E\in\tilde\Gamma^{[1]}\setminus(\tilde\Gamma^{[1]}_\infty\cup S)}M/\Z u_{(\partial^-\hat E,\hat E)}
\]
之间诱导出一个同构。

下面我们证明 $\Phi$ 的指数与 $\Phi''$ 的指数一致。由于 $\Phi''$ 的指数就是映射 $\Phi'$ 的次数，结论随之成立。

引入简记，令
\begin{align*}
L_1&=\operatorname{Map}(\hat\Gamma^{[0]},M),\\
L_2&=\prod_{E\in\hat\Gamma^{[1]}\setminus\hat\Gamma^{[1]}_\infty}M/\Z u_{(\partial^-E,E)},\\
L_3&=\prod_{i=1}^sM/\Z u_{(\partial^-E_i,E_i)},\\
M_1&=\prod_{V\in\tilde\Gamma^{[0]}\setminus\hat\Gamma^{[0]}}M/\Z u_{(\partial^-E(V),E(V))},\\
M_2&=\prod_{E\in\tilde\Gamma^{[1]}\setminus(\tilde\Gamma^{[1]}_\infty\cup S)}M/\Z u_{(\partial^-\hat E,\hat E)},
\end{align*}
于是我们有映射
\begin{align*}
\Phi&:L_1\to L_2\oplus L_3\\
\Phi''&:L_1\oplus M_1\to L_2\oplus M_2\oplus L_3,
\end{align*}
其中 $\Phi''$ 诱导出同构 $M_1\to M_2$。

现在考虑图表
\begin{equation*}
\begin{array}{ccc}
L_1 & \xrightarrow{\ \Phi\ } & L_2\oplus L_3\\[6pt]
{\scriptstyle\Psi_1}\Big\downarrow & & \Big\downarrow{\scriptstyle\Psi_2}\\[6pt]
L_1\oplus M_1 & \xrightarrow{\ \Phi''\ } & L_2\oplus M_2\oplus L_3
\end{array}
\end{equation*}
其中
\[
\Psi_1\big((m_V)_{V\in\hat\Gamma^{[0]}}\big)=\big((m_V)_{V\in\hat\Gamma^{[0]}},(m_{\partial^-E(V)})_{V\in\tilde\Gamma^{[0]}\setminus\hat\Gamma^{[0]}}\big),
\]
而 $\Psi_2$ 是自然包含。容易验证该图表交换，故由蛇引理，$\Phi$ 与 $\Phi''$ 的余核同构。但余核的阶就是像的指数，故结论成立。
\end{proof}

给定由命题 4.22 得到的环面横截前对数曲线 $f:C\to X_0$，我们现在希望计数具有给定底层概形态射的对数态射 $f:C^\dagger\to X_0^\dagger$ 的非同构选取之数。一般地，我们在 $C$ 上放置对数标记点的自由度过大，因此我们施加

% [原文 p.151]
如下条件：$f$ 在 $X_0^\dagger\to\operatorname{Spec}k^\dagger$ 为严格处也应是严格的（见定义 3.21）。注意 $X_0^\dagger\to\operatorname{Spec}k^\dagger$ 在 $\operatorname{Sing}(X_0)\cup\partial X_0$ 之外是严格的。确实，在 $X_0\setminus(\operatorname{Sing}(X_0)\cup\partial X_0)$ 中某点的邻域内，$X_0$ 上的对数结构就是 $\cM_{X_0}=\cO_{X_0}^\times\oplus\N$。因此这一严格性约束要求：$C$ 中映入 $\operatorname{Sing}(X_0)\cup\partial X_0$ 之外的点都必须是 $C$ 的光滑点，而不是对数标记点。无论如何，由 $f:C\to X_0$ 的构造，$C$ 的二重点都映到 $X_0$ 的二重点，故由命题 4.9 的 (1)，$C^\dagger$ 的任一对数标记点都必须映入 $\partial X_0$。这是施加严格性带来的唯一实质性后果。

\begin{proposition}
设 $f:(C,x_1,\dots,x_s)\to X_0$ 是由简单热带曲线 $h$ 用命题 4.22 构造出的环面横截前对数曲线。进一步假设对每个 $E\in\tilde\Gamma^{[1]}\setminus\tilde\Gamma^{[1]}_\infty$，边 $h(E)$ 的仿射长度可被 $w(E)$ 整除。（注意这总可以通过重新缩放 $M$ 来实现。）则底层概形态射为 $f$、且在 $X_0^\dagger\to\operatorname{Spec}k^\dagger$ 为严格处严格的不同构映射
\[
f^\dagger:(C^\dagger,x_1,\dots,x_s)\to X_0^\dagger
\]
的个数为
\[
\Big(\prod_{E\in\hat\Gamma^{[1]}\setminus\hat\Gamma^{[1]}_\infty}w(E)\Big)\cdot\prod_{i=1}^sw(E_i).
\]
\end{proposition}

\begin{proof}[证明]
令 $C^o\subseteq C$ 为 $C$ 的开子集
\[
C^o:=C\setminus f^{-1}\big(\operatorname{Sing}(X_0)\cup\partial X_0\big).
\]
则要求是 $f^\dagger:(C^o)^\dagger\to X_0^\dagger$ 为严格，故在 $C^o$ 上对数结构只是 $X_0$ 上对数结构的拉回。因此 $C^o$ 上的对数结构没有选择的余地。我们需要理解把这个对数结构延拓到 $C$ 上有多少种方式。

我们把讨论分为如下几部分。我们关心的是把 $f^\dagger:(C^o)^\dagger\to X_0^\dagger$ 在 $\operatorname{Spec}k^\dagger$ 上分类延拓为 $f^\dagger:C^\dagger\to X_0^\dagger$，且在同构意义下。更精确地说，设有两个这样的延拓，记作 $f_i^\dagger:C_i^\dagger\to X_0^\dagger$，$i=1,2$。这里 $C_i^\dagger$ 的底层概形都是 $C$。若存在对数同构 $\kappa^\dagger:C_1^\dagger\to C_2^\dagger$ 使 $f_1^\dagger=f_2^\dagger\circ\kappa^\dagger$，我们就说这两个延拓\emph{同构}。若这种同构 $\kappa^\dagger$ 的底层概形态射是恒等映射，我们就说它是\emph{概形意义下平凡的}。

随后我们将分三步证明该结果。第一，我们证明 $f^\dagger:(C^o)^\dagger\to X_0^\dagger$ 越过满足 $f(x)\in\partial X_0$ 的点 $x\in C$ 的延拓是唯一的。第二，我们在概形意义下平凡同构的意义下分类 $f^\dagger$ 越过 $C$ 的二重点的延拓。第三，我们计入那些不是概形意义下平凡的同构。

\emph{第 1 步. 越过满足 $f(x)\in\partial X_0$ 的点 $x\in C$ 的延拓.} 设 $x\in C$ 满足 $f(x)\in\partial X_0$。由环面横截性，$f(x)$ 落在 $X$ 的一个一维环面轨道中。设 $\sigma\in\Sigma_\cP$ 为对应于该轨道的二维锥，它必然具有形状 $\R_{\ge0}(v,1)+\R_{\ge0}(v',0)$，其中 $v\in\cP$，而 $v'\in M$ 本原、生成 $\Sigma$ 的一条射线。注意在作用 $\mathrm{GL}_3(\Z)$ 的一个元素之后，我们可以假设 $v'=(1,0)$ 且 $v=(0,0)$，于是开仿射子集 $X_\sigma\subseteq X$ 为 $X_\sigma=\operatorname{Spec}k[t,u,y^{\pm1}]$，其中 $t=z^{(0,0,1)}$ 给出正则函数 $\pi$，$u=z^{(1,0,0)}$

% [原文 p.152]
而 $y=z^{(0,1,0)}$。用这些变量，我们有 $(\partial X)\cap X_\sigma=V(tu)$。包含 $f(x)$ 的一维轨道由 $t=u=0$ 给出。设 $s_t$ 与 $s_u$ 分别表示 $\cM_X$ 在 $X_\sigma$ 上对应于 $t$ 与 $u$ 的截面。我们也用 $s_t$ 与 $s_u$ 记它们对 $\cM_{X_0}$ 的限制。于是，在 $f(x)$ 附近，$\cM_{X_0}$ 由形如 $s_t^as_u^b\varphi$ 的截面组成，其中 $\varphi$ 在 $f(x)$ 的某邻域内可逆。

我们如何在一个邻域内给 $C$ 放置对数结构，使得 $f^\dagger$ 延拓为越过 $x$ 的对数态射？假设存在 $f^\dagger:(C^o)^\dagger\to X_0^\dagger$ 的延拓 $f^\dagger:C^\dagger\to X_0^\dagger$。则在 $x$ 的某邻域内，我们必须有 $\alpha_C f^\#(s_u)=f^*(u)$，而 $f^*(u)$ 在 $x$ 的该邻域内于 $x$ 之外可逆。于是 $f^\#(s_u)$ 在 $\overline{\cM}_C$ 中的像在该邻域内恰以 $x$ 为支集。这说明 $x$ 必须是对数标记点。

有了这一点，把 $f^\dagger:(C^o)^\dagger\to X^\dagger$ 延拓到 $x$ 的方式是唯一的。在 $x$ 的邻域内，我们可以写 $\cM_C=\cM_{(C,x)}\oplus\N\rho_C$，其中 $\cM_{(C,x)}$ 是由 $x\in C$ 定义的除子型对数结构，而 $\rho_C$ 定义态射 $C^\dagger\to\operatorname{Spec}k^\dagger$。此时我们必然有 $f^\#(s_t)=\rho_C$，$f^\#(s_u)=f^*(u)\in\cM_{(C,x)}$。这就给出 $f^\dagger:(C_0)^\dagger\to X_0^\dagger$ 在 $\operatorname{Spec}k^\dagger$ 上越过 $x$ 的唯一延拓。

% [Tutor Note（原文笔误，照译未改）: 此处 $(C_0)^\dagger$ 应为 $(C^o)^\dagger$，其中 $C^o=C\setminus f^{-1}(\operatorname{Sing}(X_0)\cup\partial X_0)$；同段前文两次写作 $(C^o)^\dagger$。]

因此 $C$ 中映入 $\partial X_0$ 的点不带来任何问题：没有选择的余地。

\emph{第 2 步. 越过 $C$ 的结点 $x\in C$ 的延拓.} 设 $x\in C$ 是结点。令 $E$ 为 $\tilde\Gamma$ 中相应的边，$\ell$ 为 $h(E)$ 的仿射长度，且 $e=\ell/w(E)$。由命题 4.9 的 (4) 我们知道，所需的 $\cM_C$ 必然满足 $\cM_{C,\bar x}=S_e$。我们将证明：事实上，在概形意义下平凡同构的意义下，$x$ 的一个邻域内 $f^\dagger$ 恰有 $\mu:=w(E)$ 种选择。

把 $C$ 在 $x$ 附近局部写成 $C=\operatorname{Spec}k[z,w]/(zw)$，并把 $X$ 在 $f(x)$ 附近局部写成
\[
X=\operatorname{Spec}k[u,v,y^{\pm1},t]/(uv-t^\ell),
\]
于是
\[
X_0=\operatorname{Spec}k[u,v,y^{\pm1}]/(uv).
\]
这里，在命题 4.9 证明所用的相应记号下，$u=z^{\alpha_1}$，$v=z^{\alpha_2}$，$t=z^\rho$。注意由命题 4.9 的 (2)，在 $x$ 的一个邻域内有
\[
f^*(u)=z^\mu\cdot\mathrm{unit},\qquad f^*(v)=w^\mu\cdot\mathrm{unit}.
\]
通过在 $C$ 上 $x$ 附近作局部坐标变换，我们可以假设这些单位都是 $1$。现在 $u,v,t$ 是 $\cM_X$ 的截面；为避免混淆，我们把 $\cM_X$ 的这些截面写作 $s_u,s_v,s_t$，而把 $u,v,t$ 视为正则函数。于是我们可把 $s_u,s_v,s_t$ 限制为 $\cM_{X_0}$ 的截面，也称为 $s_u,s_v,s_t$。注意
\[
\alpha_{X_0}(s_u)=u,\qquad\alpha_{X_0}(s_v)=v,\qquad\alpha_{X_0}(s_t)=0.
\]
我们要寻找一个交换图表
\begin{equation}
\begin{array}{ccc}
f^{-1}\cM_{X_0} & \xrightarrow{\ f^\#\ } & \cM_C\\[6pt]
{\scriptstyle\alpha_{X_0}}\Big\downarrow & & \Big\downarrow{\scriptstyle\alpha_C}\\[6pt]
f^{-1}\cO_{X_0} & \xrightarrow{\ f^*\ } & \cO_C
\end{array}
\tag{4.9}
\end{equation}

% [原文 p.153]
首先，我们在 $x$ 的一个邻域内构造 $f^\dagger$ 的 $\mu$ 个延拓。对任意满足 $\zeta^\mu=1$ 的 $\zeta\in k^\times$，考虑在 $x$ 的邻域内定义的图表 $S_e\to\cO_C$，它在该邻域内定义一个对数结构 $\cM_C$，由
\begin{equation}
((a,b),c)\mapsto\begin{cases}(\zeta^{-1}z)^aw^b,&c=0,\\ 0,&c\ne0.\end{cases}
\tag{4.10}
\end{equation}
给出。这把 $\cM_C$ 局部描述为 $S_e\oplus\cO_C^\times$ 的商，故 $S_e$ 的一个元素诱导出 $\cM_C$ 的一个元素。对每个 $((a,b),c)\in S_e$，记 $s_{((a,b),c)}$ 为 $\cM_C$ 的相应截面。然后定义 $s_z=s_{((1,0),0)}$，$s_w=s_{((0,1),0)}$，使得 $\alpha_C(s_z)=\zeta^{-1}z$，$\alpha_C(s_w)=w$。我们用截面 $s_{((0,0),1)}$ 定义映射 $C^\dagger\to\operatorname{Spec}k^\dagger$；由例 3.26 的 (2)，这是对数光滑的。我们通过
\begin{equation}
f^\#(s_u)=s_z^\mu,\qquad f^\#(s_v)=s_w^\mu,\qquad f^\#(s_t)=s_{((0,0),1)}
\tag{4.11}
\end{equation}
定义 $f^\#:f^{-1}\cM_X\to\cM_C$。这使 (4.9) 交换，从而对 $\zeta$ 的每种选取，给出 $x$ 的邻域内 $C$ 上对数结构的一种选取，以及 $f^\dagger$ 越过 $x$ 的一个延拓。

% [Tutor Note: 第 2 步中图表 $S_e\to\cO_C$ 由 $\mu=w(E)$ 次单位根 $\zeta$ 决定，而 $\zeta$ 通过 $\alpha_C(s_z)=\zeta^{-1}z$ 与 $f^\#(s_t)=s_{((0,0),1)}$ 相连；$\zeta$ 的不同选取在第 3 步又被 $\operatorname{Aut}_0(C,x_1,\dots,x_s)$ 的作用（由 $\eta_i(\kappa)$ 描述）部分消解，最终恰剩 $w(E)$ 个同构类（无界情形为唯一）。]

现在有两件事需要证明。第一，我们需要证明这 $\mu$ 种选择中没有任何两种通过概形意义下平凡的同构等同起来。第二，我们需要证明所有可能的延拓都具有上述形式。

对第一点，设 $\mu$ 次根的两种选取 $\zeta_1,\zeta_2$ 给出 $C_1^\dagger,C_2^\dagger$，带有映射 $f_i^\dagger:C_i^\dagger\to X_0^\dagger$。假设存在概形意义下平凡的同构 $\kappa^\dagger:C_1^\dagger\to C_2^\dagger$，即一个交换图表
\begin{equation*}
\begin{array}{ccc}
\cM_{C_2} & \xrightarrow{\ \kappa^\#\ } & \cM_{C_1}\\[6pt]
{\scriptstyle\alpha_{C_2}}\Big\downarrow & & \Big\downarrow{\scriptstyle\alpha_{C_1}}\\[6pt]
\cO_C^\times & \xrightarrow{\ =\ } & \cO_C^\times
\end{array}
\end{equation*}
则我们必须有
\begin{equation}
\alpha_{C_2}(s_z)=\zeta_2^{-1}z=\alpha_{C_1}\big(\kappa^\#(s_z)\big)
\tag{4.12}
\end{equation}
\begin{equation}
\alpha_{C_2}(s_w)=w=\alpha_{C_1}\big(\kappa^\#(s_w)\big).
\tag{4.13}
\end{equation}
当然，由于 $\kappa$ 必须是 $\operatorname{Spec}k^\dagger$ 上的同构，我们必须有 $\kappa^\#(s_t)=s_t$。现在显然 $\kappa^\#$ 必须在 $\overline{\cM}_{C_i}$ 上诱导恒等映射，故在 $x$ 的一个邻域内必须有 $\kappa^\#(s_z)=\varphi_z\cdot s_z$ 与 $\kappa^\#(s_w)=\varphi_w\cdot s_w$，其中 $\varphi_z,\varphi_w$ 是可逆函数。但 (4.12) 告诉我们在 $z\ne0$ 的点处 $\varphi_z=\zeta_1/\zeta_2$，而 (4.13) 告诉我们在 $w\ne0$ 的点处 $\varphi_w=1$；但把 $\kappa^\#$ 作用于关系式 $s_zs_w=s_t^e$ 可得 $\varphi_z\varphi_w=1$。这只有当 $\zeta_1=\zeta_2$ 且 $\varphi_z=\varphi_w=1$ 时才可能。

对第二点，假设给定一个图表 (4.9)。考虑诱导映射 $\bar f^\#:f^{-1}\overline{\cM}_{X_0}\to\overline{\cM}_C$。在 $x$ 处的茎上，我们只能取 $\bar f^\#:S_\ell\to S_e$ 为
\begin{align*}
\bar f^\#\big((1,0),0\big)&=\big((\mu,0),0\big),\\
\bar f^\#\big((0,1),0\big)&=\big((0,\mu),0\big),\\
\bar f^\#\big((0,0),1\big)&=\big((0,0),1\big).
\end{align*}
确实，这在命题 4.9 的证明中已证：见 (4.5)。于是我们可以在 $x$ 的一个邻域内找到 $\cM_C$ 的截面 $s_z',s_w'$，使得 $f^\#(s_u)=(s_z')^\mu$，$f^\#(s_v)=(s_w')^\mu$。这些截面在相差一个 $\mu$ 次单位根的意义下是确定的。由于

% [原文 p.154]
$\alpha_C\big((s_z')^\mu\big)=z^\mu$，$\alpha_C\big((s_w')^\mu\big)=w^\mu$，在约束 $\alpha_C(s_z')=z$，$\alpha_C(s_w')=w$ 之下，$s_z',s_w'$ 的选取是唯一的。于是，把 $f^\#$ 作用于关系式 $s_us_v=s_t^\ell$ 给出 $\cM_C$ 中的 $(s_z's_w')^\mu=s_t^\ell$，即 $s_z's_w'=\zeta s_t^e$，其中 $\zeta$ 是一个 $\mu$ 次单位根。现在定义映射 $S_e\to\cM_C$ 为
\[
((a,b),c)\mapsto(\zeta^{-1}s_z')^a(s_w')^bs_t^c.
\]
把该映射与 $\alpha_C$ 复合，便给出由 (4.10) 所给对数结构 $\cM_C$ 的一个图表。在记号
\[
s_z=s_{((1,0),0)},\qquad s_w=s_{((0,1),0)}
\]
之下，我们有 $s_z'=\zeta s_z$，$s_w'=s_w$，而映射 $f^\#$ 恰如 (4.11) 所给。

\emph{第 3 步. 非平凡等同.} 到目前为止，看起来我们已经给出了 $\prod_{E\in\tilde\Gamma^{[1]}\setminus\tilde\Gamma^{[1]}_\infty}w(E)$ 种对数态射的选择。然而，我们还没有计入这些选择之间的、不是概形意义下平凡的同构。设 $f_i^\dagger:(C_i^\dagger,x_1,\dots,x_s)\to X_0^\dagger$ 是上面构造的两种选择，并设 $\kappa^\dagger:C_1^\dagger\to C_2^\dagger$ 给出它们之间的一个同构。

当然，$\kappa^\dagger$ 诱导同构 $\kappa^\#:\kappa^{-1}\cM_{C_2}\to\cM_{C_1}$，故可以假设 $\kappa^\dagger$ 是严格的，完全由底层概形自同构 $\kappa:C\to C$ 决定。

现在，由于这样的自同构必须满足 $f\circ\kappa=f$，$\kappa$ 能做的事情受到一定限制。$C$ 的自同构诱导 $\tilde\Gamma$ 的自同构 $\tilde\kappa$，由 $\tilde\kappa(V)=V'$（若 $\kappa(C_V)=C_{V'}$）定义。我们必须有 $h\circ\tilde\kappa=h$。由于 $h$ 在 $\hat\Gamma$ 的顶点集上是单射（因为 $h$ 是单的），由此可知 $\tilde\kappa$ 是恒等映射，从而对 $C$ 的每个不可约分量 $C'$ 有 $\kappa(C')=C'$。

我们来考虑保持分量的 $\operatorname{Aut}(C,x_1,\dots,x_s)$ 的元素。这些元素构成 $\operatorname{Aut}(C,x_1,\dots,x_s)$ 的单位元连通分支子群，我们记作 $\operatorname{Aut}_0(C,x_1,\dots,x_s)$。注意 $C$ 的所有分量都至少有 $3$ 个特殊点（结点、标记点或对数标记点），只有对应于 $\tilde\Gamma$ 的二价顶点的那些分量例外，它们有两个特殊点。有超过 $2$ 个特殊点的分量的自同构群是平凡的，有两个特殊点的分量的自同构群是 $\mathbb{G}_m$。因此 $\operatorname{Aut}_0(C,x_1,\dots,x_s)$ 是 $(\mathbb{G}_m)^{\#\tilde\Gamma^{[0]}_b}$，其中 $\tilde\Gamma^{[0]}_b$ 表示 $\tilde\Gamma$ 的二价顶点之集。

现在我们分析这些自同构对对数结构的作用。设 $E\in\hat\Gamma^{[1]}$，$\mu=w(E)$。先假设对任意 $i$ 都有 $E\ne E_i$。在 $\tilde\Gamma$ 中，我们可以假设 $E$ 分裂为边 $F_1,\dots,F_r$，若 $E$ 有界则 $\partial F_i=\{V_{i-1},V_i\}$（$1\le i\le r$），若 $E$ 无界则 $2\le i\le r$。若 $E$ 无界，我们还有 $\partial F_1=\{V_1\}$，于是 $F_1$ 无界。令 $p_i\in C$ 为对应于 $F_i$ 的特殊点：若 $F_i$ 有界，则 $p_i$ 是二重点；若 $F_i$ 无界，则 $f(p_i)\in\partial X_0$ 且 $p_i$ 是对数标记点。

若 $F_i$ 是有界边，则注意 $(C_{V_{i-1}}\setminus\{p_{i-1}\})\cup(C_{V_i}\setminus\{p_{i+1}\})\subseteq C$ 同构于 $\operatorname{Spec}k[z_i,w_i]/(z_iw_i)$，其中 $z_i$ 在 $C_{V_{i-1}}$ 上非零，$w_i$ 在 $C_{V_i}$ 上非零。若 $F_1$ 无界，则 $C_{V_1}\setminus\{p_2\}=\operatorname{Spec}k[w_1]$。由于每个 $C_{V_i}$ 都是 $\PP^1$，当 $F_i$ 与 $F_{i+1}$ 都有界时我们还有 $z_{i+1}=w_i^{-1}$。对 $F_i$ 有界的情形，$h(F_i)$ 是 $\cP$ 中长度为 $\ell_i$ 的边 $\omega_i$，且 $X_{C(\omega_i)}\subseteq X$ 可写成 $\operatorname{Spec}k[u_i,v_i,y_i^{\pm1},t]/(u_iv_i-t^{\ell_i})$。当 $F_1$ 无界时，$h(F_1)=\omega_1$，我们写 $X_{C(\omega_1)}=\operatorname{Spec}k[t,v_1,y_1^{\pm1}]$。在 $X_{C(\omega_i)}\cap X_{C(\omega_{i+1})}$ 上有 $u_{i+1}=v_i^{-1}$。由

% [原文 p.155]
命题 4.17 中关于二价直线的描述，注意 $f^*(u_i)=z_i^\mu$（$2\le i\le r$）且 $f^*(v_i)=w_i^\mu$（$1\le i\le r-1$）。

设 $\eta_i:\operatorname{Aut}_0(C,x_1,\dots,x_s)\to k^\times$ 为到对应于 $C_{V_i}$ 的分量的投影。则对 $\kappa\in\operatorname{Aut}_0(C,x_1,\dots,x_s)$，我们可以如下描述 $\kappa$ 在分量 $C_{V_i}$ 上的作用，该分量有坐标 $w_i=z_{i+1}^{-1}$：
\[
\kappa^*w_i=\eta_i(\kappa)\cdot w_i,\qquad\kappa^*z_{i+1}=\eta_i(\kappa)^{-1}\cdot z_{i+1}.
\]
于是用 $\kappa$ 作拉回与 $f^*u_i=z_i^\mu$，$f^*v_i=w_i^\mu$ 相容，当且仅当对所有 $i$ 有 $\eta_i(\kappa)^\mu=1$。现在考虑决定二重点 $p_1,\dots,p_r$（若 $F_1$ 无界则为 $p_2,\dots,p_r$）邻域内 $\cM_C$ 的图表的 $\mu$ 次单位根数据 $\zeta_1,\dots,\zeta_r$（或 $\zeta_2,\dots,\zeta_r$）。于是 $p_i$ 处的图表由
\[
((a,b),c)\mapsto\begin{cases}(\zeta_i^{-1}z_i)^aw_i^b,&c=0\\ 0,&c\ne0\end{cases}
\]
给出；而施加 $\kappa$ 把该图表换为图表
\[
((a,b),c)\mapsto\begin{cases}(\zeta_i^{-1}\eta_{i-1}(\kappa)^{-1}z_i)^a(\eta_i(\kappa)w_i)^b,&c=0\\ 0,&c\ne0\end{cases}
\]
该图表等价（见定义 3.18）于图表
\[
((a,b),c)\mapsto\begin{cases}\big((\eta_{i-1}(\kappa)\eta_i(\kappa)^{-1}\zeta_i)^{-1}z_i\big)^aw_i^b,&c=0\\ 0,&c\ne0\end{cases}
\]
只需把第一个图表乘以由
\[
((a,b),c)\mapsto\eta_i(\kappa)^a/\eta_i(\kappa)^b
\]
给出的函数 $S_e\to\cO_C^\times$。于是，取
\[
\eta_i(\kappa)=\prod_{j=i+1}^r\zeta_j^{-1}
\]
（$1\le i\le r-1$），我们看到施加 $\kappa$ 把 $\zeta_i$ 换成
\[
\Big(\prod_{j=i}^r\zeta_j^{-1}\Big)\Big(\prod_{j=i+1}^r\zeta_j\Big)\zeta_i=1
\]
（$2\le i\le r$）。若 $F_1$ 有界，则 $\zeta_1$ 被换成 $\prod_{i=1}^r\zeta_i$，而这是一个固定的 $\mu$ 次单位根；若 $F_1$ 无界，则我们已经消去了所有 $\zeta_i$。因此，在同构意义下，有界情形下对应于 $F_1,\dots,F_r$ 的二重点处对数结构只有 $w(E)$ 种选择，而无界情形下对应于 $F_2,\dots,F_r$ 的二重点处对数结构只有一种选择。

其次，假设 $E=E_i$。沿用上一情形的同一记号，这意味着存在某个 $1\le i\le r-1$，使 $V_i$ 是 $E_{p_i}$ 在 $\tilde\Gamma$ 中的顶点。于是 $C_{V_i}$ 有三个特殊点，故我们只能有满足 $\eta_i(\kappa)=1$ 的自同构 $\kappa$。由无标记情形的论证，此时显然有界情形下有 $w(E_i)^2$ 种可能，无界情形下有 $w(E_i)$ 种可能。这就计入了映射 $f^\dagger:(C^\dagger,x_1,\dots,x_s)\to X_0^\dagger$ 的同构类的总数。
\end{proof}

% [原文 p.156]
\begin{proof}[定理 4.14 的证明]我们只需证明恒等式
\begin{equation}
\mathfrak{D}\cdot\Big(\prod_{E\in\hat\Gamma^{[1]}\setminus\hat\Gamma^{[1]}_\infty}w(E)\Big)\cdot\Big(\prod_{i=1}^sw(E_i)\Big)=\operatorname{Mult}(h).
\tag{4.14}
\end{equation}
首先注意，$\operatorname{Mult}(h)$ 可以用 $\Gamma$ 或 $\hat\Gamma$ 来计算：$\Gamma$ 中删去一条标记无界边后变成二价的顶点对 $h$ 的重数没有贡献。我们对 $\hat\Gamma$ 的顶点数作归纳。首先，若 $\hat\Gamma$ 只有一个顶点，则 $\hat\Gamma$ 没有紧边、有三条无界边，其中两条 $E_1,E_2$ 是标记边，设其切方向为 $u_1,u_2\in M$ 本原。则 $\Phi$ 是显然的投影
\[
\Phi:M\to M/\Z u_1\times M/\Z u_2;
\]
该投影的指数恰为 $|u_1\wedge u_2|$，故
\[
\mathfrak{D}\cdot\prod_{i=1}^2w(E_i)=w_1w_2|u_1\wedge u_2|=\operatorname{Mult}(h).
\]
在一般情形，设恰有 $s$ 个标记点，则 $\hat\Gamma$ 有 $s+1=|\Delta|$ 条无界边、$s-1$ 个顶点与 $s-2$ 条有界边。至少有一条无界边 $E$ 未被标记。令 $V=\partial^-E$。若我们从 $\hat\Gamma$ 中删去边 $E$ 与顶点 $V$，则得到两个连通分支 $\hat\Gamma_1,\hat\Gamma_2$，它们各带一条新的非紧边（这两条边原先以 $V$ 为顶点）。有两种情形。

\emph{情形 1.} $\hat\Gamma_1$ 与 $\hat\Gamma_2$ 都有顶点。此时我们得到热带曲线 $h_i:\hat\Gamma_i\to M_\R$，$i=1,2$，它由把 $h$ 限制到 $\hat\Gamma_i$ 并延长新的非紧边使 $h_i$ 固有而得到。令 $E'$ 与 $E''$ 分别为 $\hat\Gamma_1$ 与 $\hat\Gamma_2$ 的这两条新的非紧边。注意 $\hat\Gamma_1$ 与 $\hat\Gamma_2$ 现在经过点 $P_1,\dots,P_s$ 的某个子集；事实上，它们必须这样分裂，使得 $h_1$ 与 $h_2$ 都不能在保持该关联性质的前提下形变；否则 $h$ 本身会沿单参数族移动。因此我们可以对 $h_1$ 与 $h_2$ 归纳地应用该结果。

令 $u_1,\dots,u_{s-2}$ 为 $\hat\Gamma$ 的有界边在 $h$ 下的像的本原切向量，其编号使得 $u_1,\dots,u_{\ell-2}$ 对应于 $\hat\Gamma_1$ 的有界边，$u_{\ell-1}$ 对应 $E'$，$u_\ell$ 对应 $E''$，而 $u_{\ell+1},\dots,u_{s-2}$ 对应于 $\hat\Gamma_2$ 的有界边。令 $v_1,\dots,v_\ell$ 为对应于 $\hat\Gamma_1$ 的标记边的本原切向量，而 $v_{\ell+1},\dots,v_s\in M$ 为 $\hat\Gamma_2$ 的相应向量。令 $\Phi_1,\Phi_2$ 分别为 (4.6) 中为 $h_1$ 与 $h_2$ 所定义的映射，$\mathfrak{D}_i$ 为 $\Phi_i$ 的余核的阶。注意我们有
\[
\Phi_1:\operatorname{Map}(\hat\Gamma_1^{[0]},M)\to\prod_{i=1}^{\ell-2}M/\Z u_i\times\prod_{j=1}^\ell M/\Z v_j=:B'
\]
以及
\[
\Phi_2:\operatorname{Map}(\hat\Gamma_2^{[0]},M)\to\prod_{i=\ell+1}^{s-2}M/\Z u_i\times\prod_{j=\ell+1}^sM/\Z v_j=:B'',
\]
而原来的映射 $\Phi$ 由
\[
\Phi:\operatorname{Map}(\hat\Gamma_1^{[0]},M)\times\operatorname{Map}(\hat\Gamma_2^{[0]},M)\times\operatorname{Map}(\{V\},M)\to B'\times B''\times M/\Z u_{\ell-1}\times M/\Z u_\ell
\]
给出，具体为
\[
(H_1,H_2,H')\mapsto\big(\Phi_1(H_1),\Phi_2(H_2),H_1(V')-H'(V),H_2(V'')-H'(V)\big),
\]

% [原文 p.157]
其中 $V'\in\hat\Gamma_1^{[0]}$，$V''\in\hat\Gamma_2^{[0]}$ 分别是与 $E'$ 和 $E''$ 相邻的顶点。则
\[
|\operatorname{coker}\Phi|=|\operatorname{coker}\Phi_1|\cdot|\operatorname{coker}\Phi_2|\cdot\big|\operatorname{coker}(M\to M/\Z u_{\ell-1}\times M/\Z u_\ell)\big|
\]
从而
\begin{align*}
&\mathfrak{D}\cdot\Big(\prod_{E\in\hat\Gamma^{[1]}\setminus\hat\Gamma^{[1]}_\infty}w(E)\Big)\cdot\Big(\prod_{j=1}^sw(E_j)\Big)\\
&\quad=\prod_{i=1}^2\Bigg(\mathfrak{D}_i\cdot\Big(\prod_{E\in\hat\Gamma_i^{[1]}\setminus\hat\Gamma_{i,\infty}^{[1]}}w(E)\Big)\Big(\prod_{E_j\in\hat\Gamma_i^{[1]}}w(E_j)\Big)\Bigg)\cdot w(E')w(E'')|u_{\ell-1}\wedge u_\ell|\\
&\quad=\operatorname{Mult}(h_1)\operatorname{Mult}(h_2)\operatorname{Mult}_V(h)\\
&\quad=\operatorname{Mult}(h),
\end{align*}
倒数第二行用到了归纳假设。

\emph{情形 2.} $\hat\Gamma_1,\hat\Gamma_2$ 之一只由一条无界边组成，设为 $\hat\Gamma_2$。此时 $\hat\Gamma_2$ 必须经过一个标记点；否则 $h$ 会沿单参数族移动。于是 $\Phi$ 取如下形式（沿用情形 1 的记号）：
\begin{align*}
\Phi&:\operatorname{Map}(\hat\Gamma_1^{[0]},M)\times\operatorname{Map}(\{V\},M)\to B'\times M/\Z u_{s-2}\times M/\Z v_s\\
&(H_1,H')\mapsto\big(\Phi_1(H_1),H_1(V')-H'(V),H'(V)\big),
\end{align*}
并且容易看出
\[
|\operatorname{coker}\Phi|=|\operatorname{coker}\Phi_1|\cdot|u_{s-2}\wedge v_s|.
\]
由此类似地得到所需结果。

% [Tutor Note: 式 (4.14) 的证明是对 $\hat\Gamma$ 顶点数的归纳。情形 1 把 $\Phi$ 分解为 $\Phi_1$、$\Phi_2$ 与投影 $M\to M/\Z u_{\ell-1}\times M/\Z u_\ell$ 的乘积，故余核阶相乘；编号约定为：$\hat\Gamma_1$ 的有界边对应 $u_1,\dots,u_{\ell-2}$，两条新边对应 $u_{\ell-1}$（$E$ 撇）与 $u_\ell$（$E$ 两撇），$\hat\Gamma_2$ 的有界边对应 $u_{\ell+1},\dots,u_{s-2}$。]
\end{proof}

\subsection{经典世界 $\to$ 对数世界}

直观上，本节要点如下。照例给定环面簇 $X_\Sigma$ 的一个退化 $\pi:X\to\mathbb{A}^1_k$。假设给定一般纤维中的一族曲线。我们希望知道，在极限下这族曲线产生 $X_0$ 上的一条环面横截对数曲线。对一般的退化选取或一般的曲线族，这未必成立。但我们将看到，总存在一种好的环面退化选取，它能接近这一目标。再加上对曲线族的一个额外限制，就能达到在 $X_0$ 上得到环面横截对数曲线的目标。

我们开始把这一点说得更精确些，并勾勒我们接下来要做的事。我们关心如下情形。照例给定由 $M_\R$ 的某个多面体分解 $\cP$ 定义的 $\pi:X\to\mathbb{A}^1_k$；一般地，我们只假设 $\cP$ 满足定义 4.5 中的条件 (1) 与 (2)，因为本节不选取点 $P_1,\dots,P_s\in M_\R$。此外假设给定一个离散赋值环 $R$，其剩余域为 $k$、分式域为 $L$，并给定一个支配映射 $\operatorname{Spec} R\to\mathbb{A}^1_k$，它把 $\operatorname{Spec} R$ 的闭点 $p$ 映到 $0\in\mathbb{A}^1_k$。我们现在想把 $X\to\mathbb{A}^1_k$ 的一般纤维中的一族曲线看成

% [原文 p.158]
一个交换图表
\begin{equation}
\begin{array}{ccc}
(C^*,x_1^*,\dots,x_k^*) & \xrightarrow{\ f^*\ } & X\setminus X_0\\[6pt]
\Big\downarrow & & \Big\downarrow{\scriptstyle\pi}\\[6pt]
\operatorname{Spec} L & \xrightarrow{\ \psi\ } & \mathbb{A}^1_k\setminus\{0\}
\end{array}
\tag{4.15}
\end{equation}
其中 $f^*$ 是环面横截稳定映射。

我们希望以某种好的方式把族 $C^*\to\operatorname{Spec} L$ 补全。一般地我们做不到这一点，但按稳定映射的稳定约化定理（命题 2.4），在基变换之后可以做到。具体地说，存在某个 $d\ge1$ 具有如下性质：若 $R'=R[t]/(t^d-u)$，其中 $u$ 是 $R$ 的一个单值化参数，则 $R'$ 是离散赋值环，其分式域 $L'=L[t]/(t^d-u)$ 是 $L$ 的 $d$ 次扩张。于是，作基变换 $(C')^*=C^*\times_LL'$，我们得到图表
\begin{equation*}
\begin{array}{ccc}
((C')^*,x_1^*,\dots,x_k^*) & \xrightarrow{\ f^*\ } & X\setminus X_0\\[6pt]
\Big\downarrow & & \Big\downarrow{\scriptstyle\pi}\\[6pt]
\operatorname{Spec} L' & \longrightarrow & \mathbb{A}^1_k\setminus\{0\}
\end{array}
\end{equation*}
该图表随后可延拓为图表
\begin{equation*}
\begin{array}{ccc}
(C',x_1,\dots,x_k) & \xrightarrow{\ f\ } & X\\[6pt]
\Big\downarrow & & \Big\downarrow{\scriptstyle\pi}\\[6pt]
\operatorname{Spec} R' & \longrightarrow & \mathbb{A}^1_k
\end{array}
\end{equation*}
其中 $f$ 现在是稳定映射。麻烦在于：即使 $f^*$ 是环面横截的，也没有理由指望 $f_0:C_p'\to X_0$ 是环面横截的：$f_0$ 可能把某些分量收缩到 $X_0$ 的一维层，或把 $C_p'$ 的某些点映到零维层。为保证这不会发生，我们需要爆破 $X$。我们将通过环面爆破来做这件事，即通过加细定义环面簇 $X$ 的扇 $\Sigma_\cP$。由于 $f$ 在 $X_0$ 之外已经性态良好，我们可以做到这一点而不改变 $X\setminus X_0$。

这样，我们就有可能把 $X$ 爆破为某个在作稳定约化之前选定的环面簇 $\tilde X$，得到图表
\begin{equation}
\begin{array}{ccc}
(C',x_1,\dots,x_k) & \xrightarrow{\ f\ } & \tilde X\\[6pt]
\Big\downarrow & & \Big\downarrow{\scriptstyle\pi}\\[6pt]
\operatorname{Spec} R' & \longrightarrow & \mathbb{A}^1_k
\end{array}
\tag{4.16}
\end{equation}
其中 $f$ 稳定，且 $f_0:C_p'\to\tilde X_0$ 现在是环面横截的：$f_0$ 的像与 $\tilde X_0$ 的零维层不相交，且 $C_p'$ 的任何不可约分量都不映入 $\tilde X_0$ 的一维层。实现这一爆破 $\tilde X$ 将占用本节的大部分篇幅。

% [Tutor Note: 作基变换 $\mathbb{A}^1_k\to\mathbb{A}^1_k,\ t\mapsto t^e$ 之后，格 $M$ 换成 $\frac1eM$（等价地，$\widetilde M$ 换成 $M\oplus e\Z$），$\widetilde{\cP}$ 的顶点才都是整点；这与第 2 章中取 $\kappa_i$ 的 $e$ 次根、把 $\widetilde{\cM}$ 换成万有覆叠的形式语言是同一件事。]

事实上，我们希望这样做使得 $\tilde X$ 来自 $\cP$ 的一个加细。这总是可以的：为得到 $\tilde X$，我们加细 $\Sigma_\cP$ 得到 $\widetilde M_\R$ 中的扇 $\tilde\Sigma$。由于我们

% [原文 p.159]
这样做时保证了 $\tilde X\setminus\tilde X_0\cong X\setminus X_0$，事实上我们只添加了由元素 $(m,r)\in\widetilde M$（$r>0$）生成的锥。于是，取
\[
\widetilde{\cP}=\{\sigma\cap(M_\R\oplus\{1\})\mid\sigma\in\tilde\Sigma\},
\]
我们便得到 $\cP$ 的一个加细 $\widetilde{\cP}$。指定 $\widetilde{\cP}$ 等价于指定 $\tilde\Sigma$。但 $\widetilde{\cP}$ 可能有非整顶点，因为我们为得到 $\tilde\Sigma$ 所添加的某些射线的本原生成元形如 $(m,r)$，其中 $r>1$。我们希望只与顶点为整点的 $\cP$ 的加细打交道。为修正这一点，我们可以选取正整数 $e$，使得对 $\widetilde{\cP}$ 的任一顶点 $v$，$ev\in M$。于是我们可以通过由 $t\mapsto t^e$ 给出的基变换 $\mathbb{A}^1_k\to\mathbb{A}^1_k$ 对图表 (4.16) 作基变换，把 $\tilde X$ 换成 $\tilde X\times_{\mathbb{A}^1}\mathbb{A}^1$ 的正规化。容易看出这就是以格 $M\oplus e\Z\subseteq\widetilde M_\R$（代替 $\widetilde M$）与扇 $\tilde\Sigma$ 给出的环面簇。该过程保持多面体分解 $\cP$ 与 $\widetilde{\cP}$，但把格 $M$ 换成格 $\frac{1}{e}M$，使 $\widetilde{\cP}$ 的所有顶点成为整点。一旦对映射 $\tilde X\to\mathbb{A}^1_k$ 作了这个基变换，我们就把 $\operatorname{Spec} R'$ 换成 $\operatorname{Spec} R'\times_{\mathbb{A}^1}\mathbb{A}^1$ 的一个不可约分量，把 $C'$ 换成 $C'\times_{\mathbb{A}^1}\mathbb{A}^1$ 的相应不可约分量。于是我们可以假设有一个图表 (4.16)，其中 $\tilde X$ 由 $M_\R$ 的一个顶点为整点的多面体分解 $\widetilde{\cP}$ 给出。随后便可用前几节的结果研究映射 $f_0:C_p'\to\tilde X_0$，并理解它何时给出一条环面横截稳定\emph{对数}曲线。

以上就是我们将要做的事情的概要。现在我们给出精确的陈述。

\begin{theorem}
照例设 $\pi:X\to\mathbb{A}^1_k$ 由某个满足定义 4.5 中条件 (1) 与 (2) 的多面体分解 $\cP$ 定义。设 $R$ 是剩余域为 $k$、分式域为 $L$ 的离散赋值环，并设给定一个支配映射 $\psi:\operatorname{Spec} R\to\mathbb{A}^1_k$，它把 $\operatorname{Spec} R$ 的闭点 $p$ 映到 $0\in\mathbb{A}^1_k$。再设给定一个交换图表 (4.15)，其中 $f^*$ 是环面横截稳定映射。则：

\begin{enumerate}[label=(\arabic*)]
\item 或许在作了由 $t\mapsto t^e$ 给出的基变换 $\mathbb{A}^1_k\to\mathbb{A}^1_k$、并把 $R$ 换成 $R[t]/(t^d-u)$（$u$ 是 $R$ 的单值化参数）之后，存在 $\cP$ 的一个顶点为整点的加细 $\widetilde{\cP}$，它定义 $X$ 的一个环面爆破 $\tilde X$，使得图表 (4.15) 延拓为图表
\begin{equation}
\begin{array}{ccc}
(C,x_1,\dots,x_k) & \xrightarrow{\ f\ } & \tilde X\\[6pt]
\Big\downarrow & & \Big\downarrow{\scriptstyle\pi}\\[6pt]
\operatorname{Spec} R & \xrightarrow{\ \psi\ } & \mathbb{A}^1_k
\end{array}
\tag{4.17}
\end{equation}
使得 $f$ 到 $\operatorname{Spec} R$ 的闭点上的纤维 $C_p$ 的限制 $f_0$ 是一条环面横截前对数曲线，且 $\psi^*(t)$ 是 $R$ 的单值化参数。
\item 给定 (1)，我们得到对数概形的诱导图表
\begin{equation}
\begin{array}{ccc}
(C^\dagger,x_1,\dots,x_k) & \xrightarrow{\ f^\dagger\ } & \tilde X^\dagger\\[6pt]
\Big\downarrow & & \Big\downarrow{\scriptstyle\pi}\\[6pt]
(\operatorname{Spec} R)^\dagger & \longrightarrow & (\mathbb{A}^1_k)^\dagger
\end{array}
\tag{4.18}
\end{equation}
\end{enumerate}

% [原文 p.160]
这里 (4.18) 中的对数结构是由 $f^{-1}(\partial\tilde X)\subseteq C$、$\partial\tilde X\subseteq\tilde X$、$\{p\}\subseteq\operatorname{Spec} R$ 以及 $\{0\}\subseteq\mathbb{A}^1_k$ 诱导出的除子型对数结构。若此外还有
\begin{itemize}[leftmargin=2em]
\item $C_p$ 亏格为零；
\item 与前对数曲线 $f_0$ 相伴的热带曲线是单的；
\item $f^{-1}\big(\partial(\tilde X\setminus\tilde X_0)\big)$ 是 $C\to\operatorname{Spec} R$ 的互不相交的截面之并，
\end{itemize}
则诱导的对数映射 $C_p^\dagger\to\operatorname{Spec}k^\dagger$ 是对数光滑的，特别地 $f_0^\dagger$ 给出一条环面横截对数曲线。
\end{theorem}

% [Tutor Note: 本块环境序列为定义 4.21、命题 4.22、命题 4.23、定理 4.24（自动编号须接 c04a 的 4.20 之后）。定理 4.24 的陈述横跨印刷 pp.159–160，按“开始者写完”原则由本块闭合；c04c 自 p.160 的 “Before proving this theorem” 起。]
